\documentclass[11pt]{article}

\usepackage[margin=1in]{geometry}
\usepackage{amsmath,amssymb,amsthm,mathtools}
\usepackage{enumitem}
\usepackage{booktabs}
\usepackage{tabularx}
\usepackage{microtype}
\usepackage[hidelinks]{hyperref}
\usepackage{xurl}

\newtheorem{definition}{Definition}
\newtheorem{remark}{Remark}
\newtheorem{lemma}{Lemma}

\newcommand{\TV}{\mathrm{TV}}

\title{Digest-Bound Replay Plans and State Declarations:\\
\large Equivalence, Minimality, and Change Control (Series Synthesis)}
\author{}
\date{Unified archive note v0.05 (2026-02-28; archive v132)}

\begin{document}
\maketitle

\begin{abstract}
The archive publishes brief, referee-grade anonymity papers while withholding large artifact bundles (logs, scripts, certificates) unless requested.
This is only credible if (i) the public receipt surface commits to \emph{which} replay workflow would validate the claim, and (ii) any long-lived system state assumptions are pinned in a comparable way.

This note standardizes two digest-bound adjuncts already used by the worked example:
\texttt{plan\_spec\_id} (a canonical hash of the replay-plan specification) and \texttt{state\_decl\_id} (a canonical hash of the state-surface declaration).
We then give a small, operational taxonomy for drift and equivalence:
which changes are ``publication-only'' (re-run the verifier), and which changes are ``claim-affecting'' (recertify).
\end{abstract}

\paragraph{Non-redundancy note.}
Other papers in this archive should cite this note (Synthesis~18) whenever they need to justify why \texttt{plan\_spec\_id} and \texttt{state\_decl\_id} are digest-bound and which drift classes trigger replay vs recertification; they should not restate the taxonomy.

\paragraph{Scope.}
This is an interface/change-control note.
It does \emph{not} propose new privacy mechanisms.
It exists so other papers can stay tight:
when a paper says ``replay plan $\pi$'' or ``state contract $\sigma$,'' it should mean a digest-bound object, not prose.

\paragraph{Imports.}
Identifier conventions are centralized in Synthesis~8.\cite{series:notation}

Contract-object vocabulary is imported from Synthesis~0.\cite{series:contractobjects}
Receipt line-item semantics are imported from Synthesis~12.\cite{series:receiptitems}
Trace-interface and threat-window identities are imported from Synthesis~3/4.\cite{series:traceifaces,series:threatwindows}
The motivating end-to-end instantiation is the worked example (Synthesis~17).\cite{series:workedexample}
State-dependent accounting motivation is imported from State~1.\cite{series:state1}
Recertification policy under drift is imported from Certified~C.\cite{series:recert}

\paragraph{Exports.}
(i) Minimal semantic cores for \texttt{plan\_spec} and \texttt{state\_decl},
(ii) digest-binding rules (\texttt{plan\_spec\_id}, \texttt{state\_decl\_id}), and
(iii) a drift/equivalence taxonomy that tells deployers when to republish, when to re-run verifiers, and when to recertify.

\section{Replay plans: what a receipt promises}
Many papers export a numeric budget, but the public object that makes the budget meaningful is a \emph{receipt line item}.\cite{series:receiptitems}
A line item cannot ship the full artifact bundle in a source-only release, so it must at least commit to:\
(a) \emph{which} replay workflow would validate the claim, and
(b) \emph{what} evidence that workflow expects.

\begin{definition}[Replay plan specification]
A \emph{replay plan specification} \texttt{plan\_spec} is a short, deterministic description of the replay workflow referenced by \texttt{replay\_hook}.
It should contain only semantics a third party must know to (i) request the right private artifacts and (ii) deterministically recompute the public outputs.
\end{definition}

\paragraph{Minimal semantic core.}
To minimize needless churn while remaining referee-grade, we recommend that a \texttt{plan\_spec} include exactly:
\begin{enumerate}[leftmargin=*]
\item \textbf{Owner and scope:} the paper/module that owns the replay semantics; the claim fields the plan is intended to validate (typically \texttt{exposure\_nf\_id}, \texttt{tw\_id}, and optional \texttt{state\_decl\_id}).
\item \textbf{Inputs:} a list of required artifacts by name \emph{and} a short input-contract statement (e.g., ``trace bundle includes raw transcripts + eligibility filter outputs + randomness commitments'').
\item \textbf{Determinism:} the intended interpreter/toolchain version or a verifier binary digest.
\item \textbf{Checks:} a small list of named predicates the verifier applies (e.g., conformance, sample-selection, accountant recomputation).
\item \textbf{Outputs:} the expected public outputs (e.g., \texttt{verifier\_report.json} schema hash; receipt acceptance bit; derived OINL fields).
\end{enumerate}
Avoid timestamps, hostnames, and other non-semantic metadata.
If deployments want to ship operator notes, they should be separate, non-hashed fields.

\begin{definition}[Digest-binding for replay plans]
Define
\[\texttt{plan\_spec\_id} := \texttt{sha256}(\texttt{canon}(\texttt{plan\_spec})),\]
where \texttt{canon} is a deterministic canonicalization (sorted keys, no whitespace, stable number formatting).
A receipt line item that includes \texttt{replay\_hook.plan\_id} should also include \texttt{replay\_hook.plan\_spec\_id} so ``same plan id, different replay'' cannot hide off-surface.\cite{series:receiptitems,series:workedexample}
\begin{remark}[Plan certificates and transparency logs are replay evidence]
A \texttt{plan\_spec\_id} can also bind operational commitments that live outside the core anonymity math: for example,
a schedule-only deployment can publish a signed plan certificate and an inclusion proof in a plan transparency log (PTL) and treat those objects as named inputs to the replay workflow.
Operational~B provides a concrete instantiation of this pattern.\cite{series:operationalB}
\end{remark}

\end{definition}

\begin{remark}[Why this does not overreach]
\texttt{plan\_spec\_id} does \emph{not} certify the privacy claim.
It certifies the \emph{replay promise}: what an auditor will do if they request the evidence.
This is the minimum needed for non-equivocation and mechanical diffing of receipts across releases.
\end{remark}

\section{State declarations: what state a bound assumes}
Many anonymity claims are implicitly stateful (caches, reputation, long-lived monitors, adaptive fallback gates).
When a bound depends on such state, a human label like \texttt{state\_contract\_id} is not a stable interface:
it can drift without changing the label.

\begin{definition}[State-surface declaration]
A \emph{state declaration} \texttt{state\_decl} is a signed adjunct that commits to which long-lived control-plane state surfaces are in scope for the claim.
It should be treated as part of the witness interface for stateful deployments.\cite{series:state1,series:workedexample}
\end{definition}

\paragraph{Minimal semantic core.}
A \texttt{state\_decl} should include:
\begin{enumerate}[leftmargin=*]
\item \textbf{Surface inventory:} a list of named state surfaces (e.g., ``fallback usage counters,'', ``cache eviction policy,'' ``monitoring summaries'') with explicit include/exclude flags.
\item \textbf{Retention and reset:} the reset unit (epoch, window) and retention horizon for each included surface.
\item \textbf{Read/write actors:} which roles can write the surface (router, monitor, operator) and what is considered public leakage vs. internal state.
\item \textbf{Binding to the claim window:} the intended \texttt{tw\_id} (or family) the declaration is meant to accompany.
\end{enumerate}
As with plans, avoid non-semantic metadata.

\begin{definition}[Digest-binding for state surfaces]
Define
\[\texttt{state\_decl\_id} := \texttt{sha256}(\texttt{canon}(\texttt{state\_decl\_core})),\]
where \texttt{state\_decl\_core} is the minimal semantic core and \texttt{canon} is a deterministic canonicalization.
Receipts for stateful claims should carry \texttt{state\_decl\_id}.\cite{series:receiptitems,series:workedexample}
\end{definition}

\begin{lemma}[Conditioning on a published state witness bounds joint leakage]\label{lem:tv_state_witness}
Let $x,x'$ be two secrets (e.g., two lookup targets) and let $(W,Y)$ be the pair of
(i) a \emph{state witness} $W=g(S)$ derived from internal state $S$ and
(ii) a published view $Y$ (receipt line items, exposure-normal-form projections, or replay outputs).
Write $P_x$ for the joint law of $(W,Y)$ under secret $x$ and $P_{x'}$ analogously.
Then
\[
\TV(P_x, P_{x'}) \le \TV(P_x^W, P_{x'}^W) \;+\; \sup_{w}\,\TV\!\big(P_x^{Y\mid W=w},\,P_{x'}^{Y\mid W=w}\big).
\]
In particular, if $W$ is public and independent of the secret (e.g., a public refresh timer or a public ``bucket-staleness'' class),
then the joint distinguishing advantage of $(W,Y)$ is controlled by the worst-case conditional leakage of $Y$ given $W$.
\end{lemma}

\begin{proof}
For each $w,y$,
\[
P_x(w,y)-P_{x'}(w,y)=\big(P_x(w)-P_{x'}(w)\big)P_x(y\mid w)+P_{x'}(w)\big(P_x(y\mid w)-P_{x'}(y\mid w)\big).
\]
Summing absolute values over $(w,y)$ and using $\sum_y P_x(y\mid w)=1$ gives
$\|P_x-P_{x'}\|_1 \le \|P_x^W-P_{x'}^W\|_1 + \sup_w \|P_x^{Y\mid w}-P_{x'}^{Y\mid w}\|_1$.
Divide by $2$ to obtain the claim.
\end{proof}

\begin{remark}[How to use this lemma in the series]
State \emph{declarations} pin which state surfaces are in scope.
When a proof actually conditions on a public state summary (rate-limit class, bucket-staleness class, refresh epoch),
publishing that summary as an explicit receipt line item (or as part of the exposure normal form) makes the dependence auditable,
and Lemma~\ref{lem:tv_state_witness} turns ``statefulness'' into a conditional-budget accounting step.
See the state series for examples of when omitting such witnesses breaks naive horizon composition.\cite{series:state1}
\end{remark}


\section{Equivalence and drift: an operational taxonomy}
This section is deliberately operational: it tells a deployer what action is required when an id changes.

\begin{table}[t]
\centering
\small
\begin{tabularx}{\linewidth}{@{}lX@{}}
\toprule
Change & Minimal required action \\
\midrule
\textbf{ENF-ID or TW-ID changes} & New claim object (new receipt line item and ABOM/OINL); any reuse must be explicit via lineage/coarsening.\cite{series:traceifaces,series:threatwindows}\\
\textbf{StateDecl changes} (\texttt{state\_decl\_id} changes) & Treat as claim-affecting for stateful bounds: republish receipt/ABOM; if the new surface could increase leakage, trigger recertification.\cite{series:state1,series:recert}\\
\textbf{PlanSpec changes} (\texttt{plan\_spec\_id} changes) & Always republish the receipt/ABOM to pin the new replay promise. Whether the privacy claim must be recertified depends on whether the verifier/accountant semantics changed (taxonomy below).\\
\bottomrule
\end{tabularx}
\caption{Digest-bound ids turn ``is this the same claim?'' into a mechanical diff.}
\label{tab:diff}
\end{table}

\subsection*{Plan-spec drift classes (P0/P1/P2)}
Let \texttt{plan\_spec} change from \texttt{plan\_spec\_id} from $h_0$ to $h_1$.
Classify the change by what an auditor would compute on the same artifact bundle:
\begin{itemize}[leftmargin=*]
\item \textbf{P0 (non-semantic):} the verifier output and acceptance predicate are unchanged (e.g., refactoring, documentation, performance).
\emph{Action:} republish the new \texttt{plan\_spec\_id}; no new privacy evidence is required.
\item \textbf{P1 (checker-strengthening / determinism fixes):} acceptance may change only by rejecting previously-accepted bundles (more checks, tighter conformance, improved determinism).
\emph{Action:} republish and re-run the verifier on the old artifact bundle; if it still accepts, the claim can be reused.
\item \textbf{P2 (claim-affecting semantics):} the accountant, sampling rule, or acceptance predicate changes in a way that could accept bundles that would previously be rejected, or could change the computed bound.
\emph{Action:} treat as recertification-relevant drift (Certified~C): mint a new claim or run incremental/targeted audits as appropriate.\cite{series:recert}
\end{itemize}
This classification is intentionally conservative: without publishing full artifacts, the public guarantee is only as stable as the published replay semantics.

\subsection*{State-decl drift classes (S0/S1/S2)}
Similarly, classify \texttt{state\_decl} drift by whether the new declaration expands what can depend on the secret:
\begin{itemize}[leftmargin=*]
\item \textbf{S0 (metadata-only):} renaming or rewording with no surface change.
\item \textbf{S1 (surface tightening):} removing a surface from scope or reducing retention.
\item \textbf{S2 (surface expansion):} adding a surface to scope, increasing retention, or allowing new actors to write secret-dependent state.
\end{itemize}
\emph{Action:} S1 can often reuse the old bound conservatively; S2 should be treated as claim-affecting unless the owning state paper supplies an explicit attenuation lemma.

\section{Checklist for keeping papers non-redundant}
A mechanism paper should not define replay/state semantics in prose.
Instead it should:
\begin{enumerate}[leftmargin=*]
\item Export a receipt line item with \texttt{exposure\_nf\_id}, \texttt{tw\_id}, and, when relevant, \texttt{state\_decl\_id}.
\item Point replay mechanics at \texttt{plan\_id} \emph{and} \texttt{plan\_spec\_id}.
\item When describing drift/updates, cite this note for the P0/P1/P2 and S0/S1/S2 rules rather than restating them.
\end{enumerate}

\begin{thebibliography}{99}\small

\bibitem{series:contractobjects}
Anonymous.
\newblock Contract Objects for Anonymous-DHT Privacy Budgets and Claims.
\newblock Synthesis~0, The-One-True-Archive (2026).

\bibitem{series:receiptitems}
Anonymous.
\newblock Receipt Line Items for Anonymity Claims: A Minimal Tuple Schema for Published Budgets.
\newblock Synthesis~12, The-One-True-Archive (2026).

\bibitem{series:traceifaces}
Anonymous.
\newblock Trace Interfaces for Anonymous DHTs.
\newblock Synthesis~3, The-One-True-Archive (2026).

\bibitem{series:threatwindows}
Anonymous.
\newblock Threat Models and Repetition Windows for Anonymous DHTs.
\newblock Synthesis~4, The-One-True-Archive (2026).

\bibitem{series:workedexample}
Anonymous.
\newblock Worked Example: Receipt/Interlock and Drift Comparison.
\newblock Synthesis~17, The-One-True-Archive (2026).

\bibitem{series:operationalB}
Anonymous.
\newblock Auditable Trace Privacy with Abort: Disclosure Receipts for Side-Channel Resistant Transport.
\newblock Operational series Paper~B, 2026. (This archive.)

\bibitem{series:state1}
Anonymous.
\newblock State-Dependent Anonymity Budgets for Anonymous DHTs.
\newblock AnonDHT State series Paper~1, The-One-True-Archive (2026).

\bibitem{series:recert}
Anonymous.
\newblock Disagreement-Gated Recertification and Boundary Audits for Anonymous DHT Routers.
\newblock Certified series Paper~C, The-One-True-Archive (2026).

\bibitem{series:notation}
Anonymous.
\newblock Notation and Minimal Model for the One-True-Archive.
\newblock Series synthesis note (Synthesis~8), The-One-True-Archive (2026).

\end{thebibliography}

\end{document}